\documentclass[11pt,a4paper]{article}
\usepackage[margin=2.2cm]{geometry}
\usepackage[utf8]{inputenc}\usepackage[T1]{fontenc}
\usepackage{amsmath,amssymb}
\usepackage{fancyhdr}\usepackage{enumitem}\usepackage{xcolor}
\usepackage{hyperref}
\pagestyle{fancy}\fancyhf{}
\rhead{18.06 Linear Algebra (MIT OCW, Spring 2020)}
\lhead{Songpengfei}
\rfoot{Page \thepage}
\newcommand{\solvertag}[1]{{\footnotesize\color{gray}\textsf{[#1]}}}

\title{Problem Set 1 Solutions}
\author{Songpengfei}
\date{2026-10-09}
\begin{document}
\maketitle

\noindent\textbf{Solved: 5/5 (100.0\%)}\quad\textit{Auto-generated by C4C homework-solver.}

\section*{Problem P1}
Warmup questions (20 points, -4 for wrong answers until 0 is reached).
The following are supposed to be familiar from 18.02, but may not be. We recommend
reading all of Strang chapter 1, but especially 1.2 for dot products. Also there is 18.02
week 1 course notes on OCW.

\paragraph{Solution.} \solvertag{mixed} 

\paragraph{(a)} $1)\ \text{unit vector}$

\paragraph{(b)} $2)\ \text{zero vector}$

\paragraph{(c)} $a, b, c \text{ 都正确（perpendicular, orthogonal, at right angles）}$

\paragraph{(d)} $(a)\ 1)\ \text{unit vector};\ (b)\ 2)\ \text{zero vector};\ (c)\ a),b),c);\ (d2)\ \text{True};\ (e)\ 2)\ 45^\circ;\ (f)\ \det A=\pm 2$

\paragraph{(d2)} $\text{True}$

\paragraph{(e)} $45^\circ\text{（选项 2）}$

\paragraph{(f)} $\det A = \pm 2$

\noindent\textbf{Answer:}\quad $\text{(a)}\ 1)\ \text{unit vector}\quad \text{(b)}\ 2)\ \text{zero vector}\quad \text{(c)}\ a, b, c \text{ 都正确（perpendicular, orthogonal, at right angles）}\quad \text{(d)}\ (a)\ 1)\ \text{unit vector};\ (b)\ 2)\ \text{zero vector};\ (c)\ a),b),c);\ (d2)\ \text{True};\ (e)\ 2)\ 45^\circ;\ (f)\ \det A=\pm 2\quad \text{(d2)}\ \text{True}\quad \text{(e)}\ 45^\circ\text{（选项 2）}\quad \text{(f)}\ \det A = \pm 2$

\medskip

\section*{Problem P2}
Warmup questions (20 points, -4 for wrong answers until 0 is reached).
Another 18.02 concept is gradients as in the week 5 18.02 course notes on OCW.

\paragraph{Solution.} \solvertag{mixed} 

\paragraph{(a)} $3)\ \text{plane not through }0$

\paragraph{(b)} $(2,3,4)$

\paragraph{(c)} $d) \text{sphere}$

\paragraph{(d)} $(x,y,z)$

\paragraph{(e)} $\left[\begin{matrix}2\\3\\4\end{matrix}\right]$

\paragraph{(f)} $\nabla f = (2x, 2y, 2z)$

\paragraph{(g)} $\text{半空间 } 2x+3y+4z<2020$

\noindent\textbf{Answer:}\quad $\text{(a)}\ 3)\ \text{plane not through }0\quad \text{(b)}\ (2,3,4)\quad \text{(c)}\ d) \text{sphere}\quad \text{(d)}\ (x,y,z)\quad \text{(e)}\ \left[\begin{matrix}2\\3\\4\end{matrix}\right]\quad \text{(f)}\ \nabla f = (2x, 2y, 2z)\quad \text{(g)}\ \text{半空间 } 2x+3y+4z<2020$

\medskip

\section*{Problem P3}
Am I a (real)vector space? (20 pts) (A few criteria to be a vector space is that real linear
combinations make sense, and remain inside the space. In particular there should be a
zero in the space. Please answer yes or no and explain briefly if I am a vector space.
(Below you may assume the matrices have real entries.)

\paragraph{Solution.} \solvertag{mixed} 

\paragraph{(a)} $\text{Yes, } M_3(\mathbb{R}) \text{ is a real vector space.}$

\paragraph{(b)} $\text{No}$

\paragraph{(c)} $\text{Yes}$

\paragraph{(d)} $\text{Yes}$

\paragraph{(e)} $\text{Yes}$

\noindent\textbf{Answer:}\quad $\text{(a)}\ \text{Yes, } M_3(\mathbb{R}) \text{ is a real vector space.}\quad \text{(b)}\ \text{No}\quad \text{(c)}\ \text{Yes}\quad \text{(d)}\ \text{Yes}\quad \text{(e)}\ \text{Yes}$

\medskip

\section*{Problem P4}
Am I a (real) vector space (20 pts) ?

\paragraph{Solution.} \solvertag{mixed} 

\paragraph{(a)} $\text{Yes}$

\paragraph{(b)} $\text{Yes, } V=\{f\mid f(0)=0\} \text{ 是实向量空间}$

\paragraph{(c)} $\text{No}$

\paragraph{(d)} $\text{Yes, } V=\{f\mid f(17)=0\} \text{ 是实向量空间}$

\paragraph{(e)} $\text{No}$

\noindent\textbf{Answer:}\quad $\text{(a)}\ \text{Yes}\quad \text{(b)}\ \text{Yes, } V=\{f\mid f(0)=0\} \text{ 是实向量空间}\quad \text{(c)}\ \text{No}\quad \text{(d)}\ \text{Yes, } V=\{f\mid f(17)=0\} \text{ 是实向量空间}\quad \text{(e)}\ \text{No}$

\medskip

\section*{Problem P5}
Using Julia as a calculator. Include screen shots and written answers. (20 points with -2
for every wrong answer)
(If you get ERROR:UndefVarError: you may have a typo or you may need
julia> using LinearAlgebra )( In many cases, copy and paste will just work)

\paragraph{Solution.} \solvertag{mixed} 

\paragraph{(a)} $225$

\paragraph{(b)} $\|x\|_2=\sqrt{55},\quad \|y\|_2=\sqrt{979}$

\paragraph{(c)} $\cos\theta=\frac{225}{\sqrt{55}\sqrt{979}}\approx 0.969,\quad \theta\approx 0.248\text{ rad}\approx 14.2^\circ$

\paragraph{(d)} $x = A^{-1} b$

\paragraph{(e)} $b-Ax \approx \mathcal{O}(10^{-16})$

\paragraph{(f)} $(b-A*x)\cdot 2^{52}\approx \text{接近整数的数}$

\paragraph{(g)} $A^2=\begin{pmatrix}30&36&42\\66&81&96\\102&126&150\end{pmatrix},\quad A.^2=\begin{pmatrix}1&4&9\\16&25&36\\49&64&81\end{pmatrix}$

\paragraph{(h)} $x = \begin{pmatrix} -0.5 \ 0.5 \ 0.5 \end{pmatrix}$

\noindent\textbf{Answer:}\quad $\text{(a)}\ 225\quad \text{(b)}\ \|x\|_2=\sqrt{55},\quad \|y\|_2=\sqrt{979}\quad \text{(c)}\ \cos\theta=\frac{225}{\sqrt{55}\sqrt{979}}\approx 0.969,\quad \theta\approx 0.248\text{ rad}\approx 14.2^\circ\quad \text{(d)}\ x = A^{-1} b\quad \text{(e)}\ b-Ax \approx \mathcal{O}(10^{-16})\quad \text{(f)}\ (b-A*x)\cdot 2^{52}\approx \text{接近整数的数}\quad \text{(g)}\ A^2=\begin{pmatrix}30&36&42\\66&81&96\\102&126&150\end{pmatrix},\quad A.^2=\begin{pmatrix}1&4&9\\16&25&36\\49&64&81\end{pmatrix}\quad \text{(h)}\ x = \begin{pmatrix} -0.5 \ 0.5 \ 0.5 \end{pmatrix}$

\medskip

\end{document}
